% Diseño compartido de cv-es.tex y cv-en.tex. Compila con pdfLaTeX (también
% en Overleaf): Archivo es la misma letra de la web y el verde es su acento.
% Las referencias van en privado/referencias-<idioma>.tex, fuera del repo, y
% solo entran si se compila con \def\conreferencias{} (build.py --privado).

\usepackage[a4paper,left=1.8cm,right=1.8cm,top=1.4cm,bottom=1.6cm,footskip=0.8cm]{geometry}
\usepackage[T1]{fontenc}
\usepackage[lining,tabular,mainfont,semibold]{Archivo}
\usepackage{microtype}
\usepackage{xcolor}
\usepackage{enumitem}
\usepackage{needspace}
\usepackage{fancyhdr}
\usepackage{lastpage}
\usepackage[hidelinks]{hyperref}

\definecolor{acento}{HTML}{2E6B4F}
\definecolor{tinta}{HTML}{191E1B}
\definecolor{gris}{HTML}{525A54}
\definecolor{linea}{HTML}{C2C9BF}

\setlength{\parindent}{0pt}
\linespread{1.05}
\frenchspacing
\brokenpenalty=10000 % ninguna palabra partida entre dos páginas
\AtBeginDocument{\color{tinta}}

% Columna de fechas a la izquierda y texto a la derecha
\newlength{\fechas}\setlength{\fechas}{2.45cm}
\newlength{\separa}\setlength{\separa}{0.4cm}

% Pie: nombre, fecha de actualización y página
\pagestyle{fancy}
\fancyhf{}
\renewcommand{\headrulewidth}{0pt}
\fancyfoot[C]{\scriptsize\color{gris}\nombre\ · \pie\ · \thepage/\pageref*{LastPage}}

% Cabecera: nombre, una línea de presentación y los contactos
\newcommand{\cabecera}[2]{%
  {\fontsize{23}{27}\selectfont\bfseries\nombre}\par\vspace{4pt}
  {\color{gris}#1}\par\vspace{3pt}
  {\small #2}\par\vspace{7pt}
  {\color{acento}\hrule height 1.2pt}\vspace{2pt}}
\newcommand{\enlace}[2]{\href{#1}{\color{acento}#2}}

% Sección: título en versalitas verdes con una raya fina; nunca sola al pie
\newcommand{\seccion}[1]{%
  \par\Needspace{5\baselineskip}\vspace{9pt}%
  {\color{acento}\bfseries\footnotesize\textls[110]{\MakeUppercase{#1}}}\par
  \vspace{2pt}{\color{linea}\hrule height 0.5pt}\vspace{6pt}}

% Lista de entradas con la fecha en la columna izquierda
% \leavevmode antes del color, para que la fecha comparta la línea base del
% cargo; \smash, para que una fecha de dos líneas no empuje el texto
\newcommand{\fecha}[1]{\smash{\parbox[t]{\fechas}{\raggedright\small\leavevmode\color{gris}#1}}}
\newlist{entradas}{description}{1}
\setlist[entradas]{leftmargin=\dimexpr\fechas+\separa\relax, labelwidth=\fechas,
  labelsep=\separa, align=left, font=\normalfont, itemsep=5pt, topsep=0pt,
  parsep=1pt, partopsep=0pt}
% Cada entrada empieza en una página con sitio para su título y dos líneas más
\newcommand{\entrada}[2][]{\Needspace{3\baselineskip}\item[\fecha{#1}]#2}

% Cargo e institución, y un detalle opcional en gris debajo
\newcommand{\cargo}[2]{{\bfseries #1}, #2\par}
\newcommand{\detalle}[1]{{\small\color{gris}#1}\par}

% Viñetas dentro de una entrada
\newlist{puntos}{itemize}{1}
\setlist[puntos]{label={\color{acento}\textbullet}, leftmargin=1.1em,
  labelsep=0.45em, itemsep=1.5pt, topsep=2pt, parsep=0pt, before=\small}
